\chapter*{Curriculum Blueprint \& Examination Specifications}
\addcontentsline{toc}{chapter}{Curriculum Blueprint \& Examination Specifications}

\section*{Course Information}
\begin{itemize}[leftmargin=8mm]
    \item \textbf{Course Code:} CSE111
    \item \textbf{Course Title:} Fundamentals of Computing \& Information Technology
    \item \textbf{Target Audience:} First-Year Engineering Students (All Branches), Lovely Professional University
    \item \textbf{Course Credits:} 4 (L: 3, T: 1, P: 0)
    \item \textbf{Course Outcomes (COs):}
    \begin{itemize}
        \item \textbf{CO1:} Comprehend foundational computer architecture, memory hierarchy, peripheral interfaces, and data storage systems.
        \item \textbf{CO2:} Analyze operating system structures, process scheduling, memory virtualization, and the software compilation lifecycle.
        \item \textbf{CO3:} Demonstrate practical command-line proficiency, file hierarchy comprehension, and user permission management in Linux environments.
        \item \textbf{CO4:} Map industry cohort specializations, technical skill roadmaps, competitive programming frameworks, and portfolio architectures.
        \item \textbf{CO5:} Evaluate computer network topologies, protocol suites (OSI/TCP-IP), IP addressing/subnetting, and network security threat vectors.
        \item \textbf{CO6:} Apply version control workflows using Git/GitHub and leverage modern Generative AI engineering and prompt paradigms.
    \end{itemize}
\end{itemize}

\vspace{3mm}

\section*{Assessment Architecture \& Weightage Matrix}

The assessment strategy combines Continuous Assessments (CA), Mid-Term Examination, and comprehensive End-Term Examination. In this volume, all 16 papers follow the user-specified 100\% objective MCQ pattern:

\begin{center}
\small
\begin{tabularx}{\textwidth}{|l|l|c|c|c|X|}
\hline
\textbf{Assessment} & \textbf{Syllabus Scope} & \textbf{Format} & \textbf{Total Qs} & \textbf{Marks} & \textbf{Duration \& Negative Marking} \\
\hline
\textbf{CA-1} & Unit 1 \& Unit 2 & 100\% MCQ & 30 & 30 & 50 Mins, $-0.25$ per incorrect option \\
\hline
\textbf{Mid-Term} & Units 1, 2 \& 3 & 100\% MCQ & 30 & 30 & 90 Mins, $-0.25$ per incorrect option \\
\hline
\textbf{CA-2} & Unit 4 \& Unit 5 & 100\% MCQ & 30 & 30 & 50 Mins, $-0.25$ per incorrect option \\
\hline
\textbf{End-Term} & Units 1 to 6 Comprehensive & 100\% MCQ & 30 & 30 & 120 Mins, $-0.25$ per incorrect option \\
\hline
\end{tabularx}
\end{center}

\vspace{3mm}

\section*{Question Distribution across Question Papers}

\begin{itemize}[leftmargin=8mm]
    \item \textbf{CA-1 Papers (Sets A, B, C, D):} 15 questions from Unit 1 (Computer System) and 15 questions from Unit 2 (Operating System).
    \item \textbf{Mid-Term Papers (Sets A, B, C, D):} 10 questions from Unit 1 (Computer System), 10 questions from Unit 2 (Operating System), and 10 questions from Unit 3 (Linux Operating System).
    \item \textbf{CA-2 Papers (Sets A, B, C, D):} 15 questions from Unit 4 (Cohorts \& Technical Skill Sets) and 15 questions from Unit 5 (Computer Networks \& Communication).
    \item \textbf{End-Term Papers (Sets A, B, C, D):} 5 questions each across all six core modules (Units 1, 2, 3, 4, 5, and 6) providing balanced, exhaustive coverage.
\end{itemize}

\newpage
